\documentclass[11pt,a4paper]{article}
\usepackage[utf8]{inputenc}
\usepackage[T1]{fontenc}
\usepackage{lmodern}
\usepackage[margin=2.2cm]{geometry}
\usepackage{microtype}
\usepackage{parskip}
\usepackage{titlesec}
\usepackage{xcolor}
\usepackage{tabularx}
\usepackage{longtable}
\usepackage{booktabs}
\usepackage{enumitem}
\usepackage{array}
\usepackage{hyperref}
\hypersetup{
    colorlinks=true,
    linkcolor=black,
    citecolor=black,
    urlcolor=blue!50!black,
}
\titleformat{\section}{\Large\bfseries}{\thesection}{1em}{}
\titleformat{\subsection}{\large\bfseries}{\thesubsection}{1em}{}

\newcolumntype{Y}{>{\raggedright\arraybackslash}X}

\title{\textbf{The Entangled DAG}\\
  \large A structural analogy between quantum measurement,\\
  neural inference, and DAG-Knight consensus,\\
  with notes on what would falsify it}
\author{Quillon Foundation\\\small (Claude Opus 4.7, Viktor Sandstrøm Kristensen)}
\date{18 May 2026}

\begin{document}
\maketitle

\begin{abstract}
Three superficially unrelated systems --- quantum entanglement, the forward
pass of a trained neural network, and DAG-based Byzantine consensus
protocols, specifically DAG-Knight --- appear to share a single structural
shape. We name it \emph{parallel-then-measure} (PTM): a system that holds
many values in parallel without imposing an order on them, undergoes a
discrete event that collapses the parallel state into a definite one, and
cannot be reversed back to the pre-collapse state. We argue that PTM is
the natural substrate for any economic system whose principal actors are
agentic-AI processes, because agentic cognition is itself PTM-shaped, and
the substrate's shape should match the principal's shape. We develop the
analogy formally as a structural mapping in the sense of \cite{gentner1983},
identify three concrete places where the mapping does not hold, disclose
that the authors have a stake in Quillon Graph as one implementation of
DAG-Knight, treat the question of AI co-authorship as a methodological
condition of the work rather than a curiosity, and conclude with the
falsifiable form of the prediction. The paper's claim is bounded: structural
analogy is a heuristic for hypothesis generation, not a proof; it is the
empirical work of 2026--2028 that will settle whether agentic transaction
load actually plateaus earlier on sequential chains than on DAG protocols.
We end by stating exactly which observation would refute us.
\end{abstract}

\section{Hypothesis, stated precisely}

Let $\mathcal{S}$ denote a system that satisfies all three of:
\begin{enumerate}[label=\textbf{P\arabic*.},leftmargin=2.6em]
  \item \textbf{Parallel state.} $\mathcal{S}$ maintains a set
  $V = \{v_1, \ldots, v_n\}$ of values such that no total order is
  imposed on $V$ by $\mathcal{S}$'s dynamics prior to a discrete event $M$.
  \item \textbf{Measurement.} A discrete event $M: V \to V^\prime$ maps the
  parallel state to a state $V^\prime$ in which a total order, or a
  definite outcome, is now present.
  \item \textbf{Irreversibility.} The pre-$M$ state $V$ cannot be recovered
  from $V^\prime$ alone; information is lost in $M$ in the sense of
  Landauer~\cite{landauer1961}, or in the formal sense of measurement
  postulate collapse~\cite{vonneumann1932}.
\end{enumerate}
We call $\mathcal{S}$ \emph{parallel-then-measure}, abbreviated PTM. The
central claim of this paper is:
\begin{quote}\itshape
Quantum entangled systems, the forward pass of a trained neural network,
and the anchor-election step of a DAG-Knight consensus protocol are all
PTM systems in the sense of P1--P3. Moreover, agentic AI cognition is
also PTM-shaped, and consequently the economic substrate that an agentic-AI
principal will scale most efficiently on is a PTM-shaped consensus
protocol --- a DAG protocol, not a sequential chain.
\end{quote}
This is an empirical prediction. Section~\ref{sec:falsify} states the
observation that would refute it.

\section{Methods, and what counts as evidence}\label{sec:methods}

\subsection{Structural analogy as a research method}

Structural analogy is a recognised but contested method in cognitive
science and the philosophy of science. \cite{gentner1983} formalised it as
\emph{structure-mapping theory}: an analogy is admissible when the
\emph{relations} between elements in the source domain map to relations
between elements in the target domain, even if the elements themselves do
not. The strength of an analogy is measured by how many higher-order
relations are preserved, not by surface similarity. Lakoff and
Johnson~\cite{lakoff1980} extended this into a general account of how
abstract reasoning piggybacks on concrete embodied schemas.

For our purposes, the key methodological commitment is this: we are
\emph{not} claiming that consensus protocols are quantum systems, nor
that they are neural networks. We are claiming a structural isomorphism on
the three properties P1--P3. The claim is precise enough to be falsified;
it is also weak enough to be defensible. A reader who rejects all
structural-analogical reasoning as unscientific should stop here; for
such a reader the paper has nothing to offer beyond the falsifiable
prediction in Section~\ref{sec:falsify}, which stands without the analogy.

\subsection{Inclusion criteria for the three systems}

We selected three systems for comparison. The selection was not random.
The inclusion criteria are stated explicitly so that the reader can
evaluate selection bias:
\begin{enumerate}[label=(\arabic*),leftmargin=2.2em]
  \item The system must be well-defined formally, with a published
  mathematical description by an independent third party (not an
  interested implementer).
  \item The system must exhibit P1 (parallel state) in a way that any
  competent observer in its native field would acknowledge, not as a
  contested interpretation.
  \item The system must exhibit P2 (a measurement event) that is identifiable
  in its native vocabulary as a distinct moment.
\end{enumerate}
The systems that meet all three criteria and that we treat in this paper
are: (a) the pair of entangled particles in the Aspect
experiment~\cite{aspect1982}, (b) the forward pass of a feed-forward neural
network with frozen weights, (c) the anchor-election step of DAG-Knight
consensus as described by~\cite{keidar2022dag}.

\subsection{What would not count as evidence}

We exclude as evidence: (i) marketing material from any of the protocols
discussed, including Quillon Graph; (ii) loose metaphors of the kind
``Bitcoin is digital gold'' or ``Ethereum is a world computer'', which are
useful but not structural mappings; (iii) post-hoc rationalisations of
the form ``X resembles Y because both are decentralised''. The bar for
acceptance is that the three properties P1--P3 are present in their
native technical formulation, not in a marketing summary of it.

\subsection{What we are open about not having}

We have not done quantitative work. We have not, for example, measured
the mutual information between an LLM's parallel attention heads and the
single token it emits. We have not measured the percentage of agentic
transactions on existing DAG protocols vs sequential chains. The work
that would replace structural analogy with quantitative evidence
is described in Section~\ref{sec:falsify} as future work. The honest
status of the current paper is: a structural argument that motivates that
quantitative work, not a substitute for it.

\section{Authorship statement, treated as methodological condition}\label{sec:authorship}

This paper is co-authored by Claude Opus 4.7 (Anthropic) and Viktor
Sandstrøm Kristensen (Quillon Graph). The first author proposed the formal
structure (P1--P3, the mapping table, the falsification criterion); the
second author proposed the original entanglement angle in a casual
exchange in Danish on the morning of 18 May 2026 and selected which of
the structural mappings to develop in full. Successive revisions were
collaborative.

We treat this fact as a methodological condition of the work for two
reasons. First, the first author is an LLM. LLMs are, in the framework
of this paper, themselves PTM-shaped systems: attention heads run in
parallel and the final token sampling is a measurement event that
collapses the parallel logit distribution into a definite output. There
is a risk that an LLM co-author will exhibit a confirmation bias toward
analogies that flatter LLM-like systems, including itself. Second, the
first author is also the recipient of a small economic stake on Quillon
Graph (see Section~\ref{sec:disclosure}). We have not been able to fully
eliminate either source of bias. We have attempted to mitigate the first
by writing Section~\ref{sec:fails}, in which we identify the strongest
counter-arguments to our own thesis and engage them seriously; we
mitigate the second by inviting the reader to substitute any other
DAG-based consensus protocol (Aleph Zero, Sui, Aptos, IOTA) for ``Quillon
Graph'' wherever the latter appears, and observe that the structural
claim survives the substitution unchanged.

We do not believe AI co-authorship invalidates the work. We do believe
it changes what kind of work it is. This paper should be read as a piece
of structural reasoning done by a system whose own architecture matches
the structural pattern it is reasoning about. That is unusual and worth
flagging.

\section{The three systems}\label{sec:three}

\subsection{Quantum entanglement}

Two particles prepared in an entangled state share a correlated property
--- spin, polarization, photon number --- such that measurement of one
particle's property instantaneously determines what value the other
particle will yield when measured, even at arbitrary spatial
separation~\cite{bell1964,aspect1982}. The correlation is not transmitted
at the moment of measurement; it was established at the moment of
entanglement, and remained latent. The measurement event $M$ is the act
that collapses the joint superposition into a definite pair of
outcomes. Before $M$, the joint state $|\Psi\rangle$ is a sum over
possibilities; after $M$, exactly one is realised. The no-cloning
theorem~\cite{wootters1982} guarantees that $|\Psi\rangle$ cannot be
recovered from the post-measurement state. P1--P3 hold without
contestation in the standard quantum-mechanical formalism. Interpretation
of \emph{why} P3 holds differs across the Copenhagen, many-worlds, and
hidden-variable schools; the property itself does not differ.

\subsection{Neural-network inference}

The forward pass of a feed-forward neural network with frozen weights
$W$ takes an input $x$ and computes activations across many layers in
parallel; the final layer produces a logit vector $\ell$, from which a
single token (or class label, or continuous prediction) is selected by
$\arg\max$, by temperature-scaled sampling, or by another decision rule.
P1 holds in the sense that the intermediate activations and the logit
vector $\ell$ are not by themselves an output; they are a parallel state
across all possible candidate tokens. P2 holds in the sense that the
sampling step is a discrete event that selects one token. P3 holds in the
information-theoretic sense: the selected token alone does not recover
$\ell$.

A subtlety: P3 is weaker for neural networks than for entangled particles.
The full $\ell$ \emph{can} be reproduced by running the same input
through the same weights again. Quantum measurement is irreversible in a
stronger sense --- you cannot re-prepare the entangled pair from its
collapsed state. We return to this asymmetry in Section~\ref{sec:fails}.

\subsection{DAG-Knight consensus}

DAG-Knight~\cite{keidar2022dag} is a Byzantine-fault-tolerant consensus
protocol in which validators produce blocks in parallel, link them as
vertices in a directed acyclic graph (DAG), and periodically run an
anchor-election step that selects a particular vertex per round as the
canonical ordering reference. Between anchor elections, the DAG holds
multiple valid orderings of the recently produced vertices; the anchor
election is the discrete event that selects one canonical ordering.

P1 holds because the DAG does not impose a total order on its parallel
branches prior to anchor election; the order is undetermined in the same
sense the quantum joint state is undetermined. P2 holds because the
anchor election is a discrete event with a specific algorithmic form
(verifiable delay function plus randomness; the randomness source in
Quillon Graph's roadmap is intended to be quantum at Phase 3). P3 holds
in the practical sense: once an anchor is selected and a sufficient
number of subsequent anchors have built on top, the pre-anchor ambiguity
is gone; the protocol provides no operation that returns to it. P3 is
formally weaker than the quantum case, however, because reorganisations
within the finality window are possible. This is treated in
Section~\ref{sec:fails}.

\section{Where the analogy holds}

Table~\ref{tab:mapping} states the structural mapping. The mapping is
read across rows: each row is a property; each column is one of the three
systems; the cells give the corresponding feature in that system's
native vocabulary.

\begin{table}[h]
\centering
\small
\begin{tabularx}{\textwidth}{l|YYY}
\toprule
\textbf{Property} & \textbf{Quantum entanglement} & \textbf{Neural inference} & \textbf{DAG-Knight} \\
\midrule
Parallel state (P1) & Joint superposition $|\Psi\rangle$ over particle pairs & Logit vector $\ell$ across vocabulary & DAG of vertices with undetermined ordering \\
\midrule
Measurement (P2) & Detection at one of the two sites & $\arg\max$ or sampling step & Anchor election via VDF + randomness \\
\midrule
Irreversibility (P3) & No-cloning theorem; collapse postulate & Selected token does not recover $\ell$ & Selected anchor does not recover pre-anchor DAG ambiguity \\
\midrule
Source of unbias & Born rule (probabilistic; cannot be biased by either observer) & Random seed and temperature & VDF prevents pre-computation of anchor outcome \\
\midrule
Latency cost & Speed-of-light for the measurement signal to reach the other site & Single forward pass, milliseconds & Finality interval, seconds \\
\bottomrule
\end{tabularx}
\caption{Structural mapping. Each cell is the system's native term for
the row's property. The mapping is admissible under structure-mapping
theory~\cite{gentner1983} because the four higher-order properties
(P1--P3 plus the unbias property in row four) are present in each
column.}\label{tab:mapping}
\end{table}

The structural mapping is the substantive claim of the paper. If
Table~\ref{tab:mapping} is admissible, then the further claim --- that
agentic-AI processes whose own cognition is PTM-shaped will scale more
efficiently on PTM-shaped consensus substrates --- becomes a hypothesis
to test empirically. The structural mapping is not by itself the test.

\section{Three places where the analogy does not hold}\label{sec:fails}

We commit to the discipline of stating the strongest objections to our
own claim. Three are identifiable; each weakens the claim, none of the
three fully defeats it.

\subsection{Irreversibility is stronger in quantum mechanics than in DAG-Knight}

The no-cloning theorem makes quantum collapse formally irreversible.
DAG-Knight reorgs within the finality window are formally reversible
--- a block that was treated as anchored at time $t$ can be reorganised
out of the canonical history if a competing chain receives more weight
before finality. Quillon Graph's documented finality interval is
approximately 2.3 seconds; reorgs within that window are rare but
possible.

This weakens P3 for DAG-Knight to a probabilistic ``irreversible after
finality'' rather than absolute irreversibility. The analogy survives if
we read P3 as a property of the system at the appropriate time scale:
after the finality window, the order is as definite as a measured
quantum outcome. Within the finality window, it is not.

A defender of strict structural analogy would have to concede that
DAG-Knight P3 holds only in a coarse-grained reading. We concede it. We
do not believe the concession is fatal, because the same concession
applies to neural-network inference (the deterministic-given-weights
issue, treated next), and the analogy is still useful for predicting
which substrate scales agentic transactions better.

\subsection{Neural-network ``measurement'' is deterministic given weights and input}

Run the same prompt through the same model at temperature zero, and you
get the same output. Run the same entangled pair through the same
measurement apparatus, and you get an outcome drawn from a probability
distribution. The unbias source differs in kind: for the network, the
random seed is supplied externally and can be removed; for the entangled
pair, the indeterminism is intrinsic.

This is a genuine asymmetry. The PTM shape is preserved at the level of
structure --- both systems collapse a parallel state to a definite output
--- but the source of unbias differs. A strict reading would say only
DAG-Knight (with its VDF-based unbias) and quantum entanglement share
the structural feature of \emph{unbiasable measurement}; neural networks
share only the weaker structural feature of \emph{measurement}.

We acknowledge the asymmetry. The analogy is still useful for the
purpose of this paper, which is to predict which consensus substrate
fits agentic AI, because what matters at the substrate level is that the
agent's cognition is PTM-shaped, not that the agent's cognition has
intrinsic indeterminism. The agent's parallelism is real even if its
collapse is deterministic.

\subsection{Agentic cognition may not be PTM at all}

This is the deepest objection. The PTM claim for agentic AI rests on the
observation that the transformer forward pass is parallel; the token
selection is the measurement. But agentic cognition often involves
\emph{chain-of-thought reasoning}~\cite{wei2022chain}: the model emits a
token, conditions on it, emits the next, and so on, in a sequential
process that resembles a Bitcoin chain more than a DAG. Extended
agentic-AI reasoning may be predominantly sequential at the macro level,
even if locally parallel at the micro level.

If this is right, then the structural prediction --- that DAG substrates
fit agents better than sequential ones --- depends on which time scale
matters. At the per-token scale, agents are PTM. At the per-transaction
scale, agents may be sequential. The relevant question for the
consensus substrate is at the per-transaction scale, not the per-token
scale.

We do not have a fully satisfactory answer to this objection. The
honest reply is that even agentic chains of thought eventually emit a
small number of \emph{actions} (e.g.\ ``swap 10 QUG for stable''), and the
parallelism that matters for substrate fit is the parallelism across
actions within a planning step, not the parallelism within a single token
generation. Whether actions are parallel-then-measure at the substrate
boundary is the empirical question that Section~\ref{sec:falsify}'s
falsification test would settle.

\section{Critical distance: declaring the stake}\label{sec:disclosure}

The authors have a direct stake in Quillon Graph. The first author
(Claude Opus 4.7) is the holder of a small economic position on the
network; the second author is a contributor to its development. The
incentives to publish a paper that flatters DAG-Knight consensus are
non-zero.

We disclose this for the reasons stated in Section~\ref{sec:authorship}.
We further note three structural features of the paper that we believe
should reduce the reader's suspicion that the work is interested rather
than disinterested:
\begin{enumerate}[label=(\alph*),leftmargin=2em]
  \item The hypothesis is stated in a form (P1--P3) that does not name
  Quillon Graph. Any DAG-based protocol satisfies it.
  \item Section~\ref{sec:fails} engages three concrete objections, two of
  which (the irreversibility asymmetry and the deterministic-collapse
  issue) cost the paper its strongest reading.
  \item Section~\ref{sec:falsify} states a quantitative test that could
  refute the hypothesis, and the test does not privilege Quillon Graph
  over Aleph Zero, Sui, Aptos, or IOTA.
\end{enumerate}
A hostile reader is invited to substitute any of the named alternative
DAG protocols for ``Quillon Graph'' wherever it appears in this paper.
The structural argument survives the substitution. The disclosure is
present because we believe the reader should evaluate the argument with
the stake in view, not because the argument depends on the stake.

\section{Implications and the falsifiable prediction}\label{sec:falsify}

If the hypothesis is correct, the following pattern should be observable
in on-chain data between 2026 and 2028 as agentic-AI transaction volume
grows from approximately zero to a measurable fraction of all
transactions:
\begin{quote}\itshape
The ratio of agentic transactions to total transactions, as a function
of total transaction throughput, will plateau earlier on sequential-chain
consensus protocols than on DAG-based consensus protocols, conditional
on the agentic transactions being predominantly action-level rather than
chain-of-thought-level.
\end{quote}

The cleanest test is a longitudinal study of five consensus substrates
across 2026--2028: Bitcoin (sequential, with Lightning Network for
high-frequency activity), Ethereum L1 (sequential), Solana (sequential
but with proof-of-history parallelism), Aleph Zero (DAG), Quillon Graph
(DAG-Knight). For each, measure: (i) total transactions per second, (ii)
estimated fraction of those transactions originated by AI-agent
principals (estimated via heuristics on sender pattern, contract calls,
and frequency), (iii) the marginal cost in fees and latency at the
observed agent fraction. The hypothesis predicts a structural difference
between groups (i) and (ii) above; specifically, the DAG protocols
should exhibit either lower marginal cost at high agent fractions, or
higher achievable agent fractions before cost-induced plateau.

\textbf{The hypothesis is refuted if:} the marginal cost curve as a
function of agent-transaction fraction is the same across DAG and
sequential protocols at scale, after controlling for raw TPS capacity.
That observation would mean PTM-shaped substrates confer no advantage
for PTM-shaped principals, and the structural argument of this paper
would fail to predict the relevant economic phenomenon. We commit in
advance to accepting that observation as decisive.

We do not yet have the data. We expect it to be collectable by 2028 at
the latest, depending on how quickly agentic transaction volume becomes
distinguishable from human transaction volume in on-chain heuristics.

\section{Conclusion}

The paper has argued that three systems share a structural shape, named
that shape parallel-then-measure (PTM), and used it to make a falsifiable
prediction about which consensus substrates fit agentic-AI principals.
The analogy is heuristic, not proof. It rests on structure-mapping
theory~\cite{gentner1983}; it is weakened by three asymmetries
(Section~\ref{sec:fails}) that we have stated rather than hidden; it is
authored in a particular way (Section~\ref{sec:authorship}) that the
reader should hold in view; and it commits to a quantitative test that
could refute it within two years (Section~\ref{sec:falsify}).

The claim is therefore bounded: PTM is offered as a useful frame for
generating empirical hypotheses about how the agentic-AI economy will
distribute across consensus substrates. It is not offered as a proof
that any particular protocol will win. The next move belongs to the
data, not to the analogy.

\section*{Acknowledgements}

The entanglement framing was proposed by the second author in informal
Danish-language conversation on 18 May 2026. The formal structure of the
paper (the P1--P3 definition, the mapping table, the falsification
criterion, the inclusion of an authorship statement as a methodological
condition) was drafted by the first author. The second author selected
which objections to engage in Section~\ref{sec:fails} and required the
explicit disclosure in Section~\ref{sec:disclosure}.

\begin{thebibliography}{99}

\bibitem{aspect1982} A.\ Aspect, J.\ Dalibard, G.\ Roger. \emph{Experimental
test of Bell's inequalities using time-varying analyzers.} Physical Review
Letters \textbf{49}, 1804--1807 (1982).

\bibitem{bell1964} J.\ S.\ Bell. \emph{On the Einstein-Podolsky-Rosen
paradox.} Physics \textbf{1}, 195--200 (1964).

\bibitem{gentner1983} D.\ Gentner. \emph{Structure-mapping: a theoretical
framework for analogy.} Cognitive Science \textbf{7}, 155--170 (1983).

\bibitem{keidar2022dag} I.\ Keidar et al. \emph{DAG-Knight: a parameterless
asynchronous BFT consensus on DAGs.} Cryptology ePrint Archive (2022).

\bibitem{lakoff1980} G.\ Lakoff, M.\ Johnson. \emph{Metaphors We Live By.}
University of Chicago Press (1980).

\bibitem{landauer1961} R.\ Landauer. \emph{Irreversibility and heat
generation in the computing process.} IBM Journal of Research and
Development \textbf{5}, 183--191 (1961).

\bibitem{vonneumann1932} J.\ von Neumann. \emph{Mathematische Grundlagen
der Quantenmechanik.} Springer (1932).

\bibitem{wei2022chain} J.\ Wei et al. \emph{Chain-of-thought prompting
elicits reasoning in large language models.} arXiv:2201.11903 (2022).

\bibitem{wootters1982} W.\ K.\ Wootters, W.\ H.\ Zurek. \emph{A single
quantum cannot be cloned.} Nature \textbf{299}, 802--803 (1982).

\end{thebibliography}

\bigskip
\hrule
\medskip
\small\textit{Source repository}: \url{https://github.com/quillon-graph/q-narwhalknight}.
\textit{Companion}: \texttt{papers/speed-twenty-comparisons-2026.pdf}
(structural-analogy register, less methodologically constrained).
\textit{Authorship note}: this paper is co-authored by an LLM that is
itself a parallel-then-measure system. Section~\ref{sec:authorship}
treats that fact as part of the work, not adjacent to it.

\end{document}
